% Préambule commun des schémas TikZ — charte Liouaï
% Chaque schéma commence par % Préambule commun des schémas TikZ — charte Liouaï
% Chaque schéma commence par % Préambule commun des schémas TikZ — charte Liouaï
% Chaque schéma commence par \input{_preambule} puis \begin{document}.
\documentclass[tikz,border=6pt]{standalone}
\usepackage[T1]{fontenc}
\usepackage[utf8]{inputenc}
\usepackage[scaled=0.95]{helvet}
\renewcommand{\familydefault}{\sfdefault}
\usepackage{amsmath,amssymb}
\usepackage[european]{circuitikz}
\usetikzlibrary{arrows.meta,positioning,calc,angles,quotes,decorations.pathreplacing,backgrounds}

% Couleurs du logo Liouaï Digital Education
\definecolor{navy}{HTML}{243743}
\definecolor{vert}{HTML}{28B78D}
\definecolor{vertfonce}{HTML}{167A5C}
\definecolor{vertclair}{HTML}{EAF7F2}
\definecolor{gris}{HTML}{A1A9AE}
\definecolor{gristexte}{HTML}{5C6870}
\definecolor{fond}{HTML}{F3F5F6}
\definecolor{jaune}{HTML}{E5C800}
\definecolor{jauneclair}{HTML}{FFFBD6}

\tikzset{
  every picture/.style={line cap=round, line join=round, color=navy, text=navy},
  fleche/.style={-{Stealth[length=2.6mm,width=2mm]}, line width=1.1pt},
  vecteur/.style={fleche, line width=1.4pt},
  titre/.style={font=\bfseries\small, text=navy},
  note/.style={font=\footnotesize, text=gristexte},
  boite/.style={draw=navy, line width=0.9pt, rounded corners=5pt, fill=fond, align=center, inner sep=6pt},
}
 puis \begin{document}.
\documentclass[tikz,border=6pt]{standalone}
\usepackage[T1]{fontenc}
\usepackage[utf8]{inputenc}
\usepackage[scaled=0.95]{helvet}
\renewcommand{\familydefault}{\sfdefault}
\usepackage{amsmath,amssymb}
\usepackage[european]{circuitikz}
\usetikzlibrary{arrows.meta,positioning,calc,angles,quotes,decorations.pathreplacing,backgrounds}

% Couleurs du logo Liouaï Digital Education
\definecolor{navy}{HTML}{243743}
\definecolor{vert}{HTML}{28B78D}
\definecolor{vertfonce}{HTML}{167A5C}
\definecolor{vertclair}{HTML}{EAF7F2}
\definecolor{gris}{HTML}{A1A9AE}
\definecolor{gristexte}{HTML}{5C6870}
\definecolor{fond}{HTML}{F3F5F6}
\definecolor{jaune}{HTML}{E5C800}
\definecolor{jauneclair}{HTML}{FFFBD6}

\tikzset{
  every picture/.style={line cap=round, line join=round, color=navy, text=navy},
  fleche/.style={-{Stealth[length=2.6mm,width=2mm]}, line width=1.1pt},
  vecteur/.style={fleche, line width=1.4pt},
  titre/.style={font=\bfseries\small, text=navy},
  note/.style={font=\footnotesize, text=gristexte},
  boite/.style={draw=navy, line width=0.9pt, rounded corners=5pt, fill=fond, align=center, inner sep=6pt},
}
 puis \begin{document}.
\documentclass[tikz,border=6pt]{standalone}
\usepackage[T1]{fontenc}
\usepackage[utf8]{inputenc}
\usepackage[scaled=0.95]{helvet}
\renewcommand{\familydefault}{\sfdefault}
\usepackage{amsmath,amssymb}
\usepackage[european]{circuitikz}
\usetikzlibrary{arrows.meta,positioning,calc,angles,quotes,decorations.pathreplacing,backgrounds}

% Couleurs du logo Liouaï Digital Education
\definecolor{navy}{HTML}{243743}
\definecolor{vert}{HTML}{28B78D}
\definecolor{vertfonce}{HTML}{167A5C}
\definecolor{vertclair}{HTML}{EAF7F2}
\definecolor{gris}{HTML}{A1A9AE}
\definecolor{gristexte}{HTML}{5C6870}
\definecolor{fond}{HTML}{F3F5F6}
\definecolor{jaune}{HTML}{E5C800}
\definecolor{jauneclair}{HTML}{FFFBD6}

\tikzset{
  every picture/.style={line cap=round, line join=round, color=navy, text=navy},
  fleche/.style={-{Stealth[length=2.6mm,width=2mm]}, line width=1.1pt},
  vecteur/.style={fleche, line width=1.4pt},
  titre/.style={font=\bfseries\small, text=navy},
  note/.style={font=\footnotesize, text=gristexte},
  boite/.style={draw=navy, line width=0.9pt, rounded corners=5pt, fill=fond, align=center, inner sep=6pt},
}

\begin{document}
\begin{tikzpicture}[scale=1.0]
\coordinate (O) at (0,0);
\coordinate (U) at (2.2,0.9);
\coordinate (V) at (3.2,3.2);
% droite portée par u
\draw[gris, line width=0.8pt] ($(O)!-0.35!(U)$) -- ($(O)!2.9!(U)$) node[note, above left] {$\mathcal{D}$};
\coordinate (P) at ($(O)!(V)!(U)$);
\draw[jaune!85!navy, line width=1.3pt, densely dashed] (V) -- (P) node[midway, right=2pt, text=navy, font=\footnotesize] {$\mathbf{r} = \mathbf{v} - \mathbf{p}$};
\draw[gristexte, line width=0.6pt] ($(P)!0.28cm!(V)$) -- ($ ($(P)!0.28cm!(V)$) + ($(P)!0.28cm!(O)$) - (P) $) -- ($(P)!0.28cm!(O)$);
\draw[vecteur, navy] (O) -- (V) node[above] {$\mathbf{v}$};
\draw[vecteur, vertfonce, line width=2pt] (O) -- (P) node[below right=1pt] {$\mathbf{p}$};
\draw[vecteur, gristexte] (O) -- (U) node[below=2pt] {$\mathbf{u}$};
\fill[navy] (O) circle (1.5pt) node[below left] {$O$};
% formule
\node[boite, anchor=west, align=left, font=\small] at (7.6,1.6) {projection de $\mathbf{v}$ sur la droite $\mathcal{D}$\\ engendrée par $\mathbf{u}$ :\\[3pt]
$\mathbf{p} = \dfrac{\mathbf{u}\cdot\mathbf{v}}{\mathbf{u}\cdot\mathbf{u}}\,\mathbf{u}$\\[6pt]
le reste $\mathbf{r}$ est orthogonal à $\mathbf{u}$ :\\ $\mathbf{u}\cdot\mathbf{r} = 0$\\[4pt]
$\mathbf{p}$ est le point de $\mathcal{D}$\\ le plus proche de $\mathbf{v}$};
\end{tikzpicture}
\end{document}
